% !TEX root = task-1-1-5.tex
% Editable task bank. Append new BankTask blocks at the end of the document.
\documentclass[12pt,a4paper]{article}
\usepackage{physics-textbook}
\usepackage{graphicx,needspace}
\usetikzlibrary{decorations.pathmorphing}
\newif\ifShowSolutions
\ifdefined\TutorEdition
  \ShowSolutionstrue
\else
  \ShowSolutionsfalse % Change to \ShowSolutionstrue to display solutions.
\fi
\newcounter{banktask}
\newsavebox{\BankTaskBox}
\newcommand{\BankTaskContent}[3]{%
  \noindent\textbf{\thebanktask.}\quad #2\par
  \ifShowSolutions
    \Needspace{50mm}% Keep the solution heading with its opening text or diagram.
    #3\par
    \ifdefined\TutorEdition
      \ifcsname TutorSource#1\endcsname
        \nopagebreak\smallskip{\small\csname TutorSource#1\endcsname\par}%
      \fi
    \fi
  \fi
}
\newcommand{\BankTask}[3]{%
  \refstepcounter{banktask}\par\medskip
  \sbox{\BankTaskBox}{\begin{minipage}{\linewidth}%
    \BankTaskContent{#1}{#2}{#3}\end{minipage}}%
  \ifdim\dimexpr\ht\BankTaskBox+\dp\BankTaskBox\relax>\dimexpr\textheight-20mm\relax
    \Needspace{35mm}\BankTaskContent{#1}{#2}{#3}%
  \else
    \Needspace{\dimexpr\ht\BankTaskBox+\dp\BankTaskBox\relax}%
    \noindent\usebox{\BankTaskBox}\par
  \fi
  \medskip
}
\begin{document}
\BookSubsection{1.1.5}{Скорость}


% Source sdamgia-3540: https://phys-ege.sdamgia.ru/problem?id=3540
\BankTask{sdamgia-3540}{
Велосипедист, двигаясь под уклон, проехал расстояние между двумя пунктами со скоростью, равной $15\,\text{км/ч}$. Обратно он ехал вдвое медленнее. Какова средняя путевая скорость на всем пути? (В километрах в час.)
}{
\smallskip\noindent\textbf{Решение}\par
Необходимо различать два понятия: среднюю путевую скорость и среднюю скорость по перемещению. Средняя путевая скорость определяется как скорость прохождения пути: $v _{\text{ср}}= \frac{S}{t} .$ То есть буквально надо весь пройденный телом путь разделить на все время, затраченное им на этот путь. Средняя путевая скорость представляет собой число, скаляр. 

Разберемся теперь со второй средней скоростью. Средняя скорость по перемещению — это вектор, равный отношению перемещения ко времени, за которое оно совершено: $\vec v _{\text{ср}}= \frac{\Delta \vec r}{\Delta t} .$ В нашей конкретной задаче, поскольку велосипедист вернулся в исходную точку, его перемещение равно нулю, а значит, его средняя скорость по перемещению тоже равна нулю.

Вычислим теперь среднюю путевую скорость. Обозначим расстояние между двумя пунктами через $L$, тогда весь путь, пройденный велосипедистом, равен $S=2L.$ На первую половину пути велосипедист затратил время $t_1= \frac{L}{v} .$ На обратную дорогу — время $t_2= \frac{L}{v /2} = \frac{2L}{v} .$ Все время пути составило $t=t_1 + t_2.$ Окончательно находим, что средняя путевая скорость велосипедиста равна: 

\[v_{\text{ср}}=\frac{S}{t}=\frac{2L}{\frac{L}{v}+\frac{2L}{v}}=\frac{2v}{3}=\frac{2\cdot15\,\text{км/ч}}{3}=10\,\text{км/ч}.\]

Ответ: $10$.
}

% Source sdamgia-6038: https://phys-ege.sdamgia.ru/problem?id=6038
\BankTask{sdamgia-6038}{
Мотоцикл едет по прямой дороге с постоянной скоростью $50\,\text{км/ч}$. По той же дороге навстречу ему едет автомобиль с постоянной скоростью $70\,\text{км/ч}$. Чему равен модуль скорости движения мотоцикла относительно автомобиля? (В километрах в час.)
}{
\smallskip\noindent\textbf{Решение}\par
Перейдем в систему отсчета, связанную с автомобилем. Модуль скорости движения мотоцикла в данной системе отсчета будет равен $50 + 70 = 120\,\text{км/ч}$.

Ответ: $120$.
}

% Source sdamgia-111: https://phys-ege.sdamgia.ru/problem?id=111
\BankTask{sdamgia-111}{
Тело разгоняется на прямолинейном участке пути, при этом зависимость пройденного телом пути $S$ от времени $t$ имеет вид:

\[S=4t + t ^2 .\]

Чему равна скорость тела в момент времени $t = 2\,\text{с}$ при таком движении? (В метрах в секунду.)
}{
\smallskip\noindent\textbf{Решение}\par
При равноускоренном движении зависимость пройденного телом пути от времени в общем виде имеет вид:

\[S= v _0 t + \frac{at ^2}{2} .\]

Сравнивая с выражением, данным в условии, заключаем, что оно укладывается в это общее правило, а значит, тело двигалось равноускоренно. Сопоставляя конкретные члены в выражениях, получаем, что начальная скорость равна $v _0=4\,\text{м/с},$ а ускорение $a=2\,\text{м/с} ^2 .$ Таким образом, скорость тела в момент времени $t=2\,\text{с}$ равна $v ( t ) = v _0 + at$:

\[v ( t=2\,\text{с} ) = v _0 + a \cdot  2\,\text{с}=4\,\text{м/с} + 2\,\text{м/с} ^2 \cdot  2\,\text{с}=8\,\text{м/с}.\]

Ответ: $8$.
}

% Source sdamgia-317: https://phys-ege.sdamgia.ru/problem?id=317
\BankTask{sdamgia-317}{
Пловец плывет по течению реки. Определите скорость пловца относительно берега, если скорость пловца относительно воды $0{,}4\,\text{м/с}$, а скорость течения реки $0{,}3\,\text{м/с}$. (В метрах в секунду.)
}{
\smallskip\noindent\textbf{Решение}\par
Вектор скорости пловца относительно берега есть сумма векторов скорости пловца относительно воды и скорости течения реки: $\vec v_{\text{пл–бер}}=\vec v_{\text{пл–вода}}+\vec v_{\text{вода}}.$ Поскольку пловец плывет по течению реки, получаем, что для величин скоростей выполняется соотношение: $v_{\text{пл–бер}}=v_{\text{пл–вода}}+v_{\text{вода}}=0{,}4\,\text{м/с}+0{,}3\,\text{м/с}=0{,}7\,\text{м/с}.$

Ответ: $0{,}7$.
}

% Source sdamgia-41130: https://phys-ege.sdamgia.ru/problem?id=41130
\BankTask{sdamgia-41130}{
\begin{center}\begin{tikzpicture}[task graph,x=.25cm,y=.025cm]
\draw[task grid,xstep=2,ystep=20](0,0)grid(18,100);
\draw[task axis](0,0)--(19.5,0)node[right]{$t,\,\text{с}$};\draw[task axis](0,0)--(0,115)node[above]{$x,\,\text{м}$};
\node[below left]at(0,0){$0$};\foreach\x in{2,4,...,18}\node[below]at(\x,0){$\x$};\foreach\y in{20,40,...,100}\node[left]at(0,\y){$\y$};
\draw[task curve](0,0)--(4,40)--(12,60)--(18,100);
\end{tikzpicture}\end{center}Спортсмен бежит по прямолинейному участку беговой дорожки. На рисунке представлен график зависимости координаты х спортсмена от времени $t$ (ось ОХ направлена вдоль беговой дорожки). Определите проекцию скорости спортсмена на ось ОХ в интервале времени от $4\,\text{с}$ до $12\,\text{с}$. (В метрах за секунду с учётом знака проекции.)
}{
\smallskip\noindent\textbf{Решение}\par
Проекция скорости равна $v _x= \frac{x_2 - x_1}{t_2 - t_1} = \frac{60 - 40}{12 - 4} = 2{,}5\,\text{м/с}.$

Ответ: $2{,}5$.
}

% Source sdamgia-41558: https://phys-ege.sdamgia.ru/problem?id=41558
\BankTask{sdamgia-41558}{
\begin{center}\begin{tikzpicture}[task graph,x=.35cm,y=.16cm]
\draw[task grid,xstep=2,ystep=5](0,0)grid(12,15);
\draw[task axis](0,0)--(14,0)node[right]{$t,\,\text{с}$};\draw[task axis](0,0)--(0,18)node[above]{$\upsilon,\,\text{м/с}$};
\node[below left]at(0,0){$0$};\foreach\x in{2,4,...,12}\node[below]at(\x,0){$\x$};\foreach\y in{5,10,15}\node[left]at(0,\y){$\y$};
\draw[task curve](0,0)--(4,15)--(12,15);
\end{tikzpicture}\end{center}На рисунке представлен график зависимости модуля скорости $\upsilon $ движущегося прямолинейно тела от времени $t$. Найдите среднюю скорость тела за время от $0$ до $10\,\text{с}$. (В метрах за секунду.)
}{
\smallskip\noindent\textbf{Решение}\par
Найдем путь, пройденный телом за $10\,\text{с}$, как площадь фигуры под графиком:

\[s= \frac{10 + 6}{2} \cdot  15=120\,\text{м}.\]

Средняя скорость за этот интервал

\[v _{\text{ср}}= \frac{s}{t} = \frac{120}{10} =12\,\text{м/с} .\]

Ответ: $12$.
}

% Source sdamgia-31875: https://phys-ege.sdamgia.ru/problem?id=31875
\BankTask{sdamgia-31875}{
\begin{center}\begin{tikzpicture}[task graph,x=.6cm,y=.6cm]
\draw[task grid,step=1](0,0)grid(6,5);
\draw[task axis](0,0)--(7,0)node[right]{$x,\,\text{м}$};\draw[task axis](0,0)--(0,5.8)node[above]{$y,\,\text{м}$};
\node[below left]at(0,0){$0$};\foreach\x in{1,...,6}\node[below]at(\x,0){$\x$};\foreach\y in{1,...,5}\node[left]at(0,\y){$\y$};
\draw[task curve,->](0,0)--(1,2)node[above left]{$\vec r_1$};\draw[task curve,->](0,0)--(5,5)node[right]{$\vec r_2$};
\end{tikzpicture}\end{center}Небольшое тело движется по плоскости YOX. На рисунке показаны радиус-векторы $\vec r_1$ и $\vec r_2$ этого тела в моменты времени $t_{1}$ = $0\,\text{с}$ и $t_{2}$ = $2{,}5\,\text{с}$. Чему равен модуль средней скорости этого тела за указанный промежуток времени? (В метрах в секунду.)
}{
\smallskip\noindent\textbf{Решение}\par
Используя данные из графика, найдем проекции радиус-вектора на координатные оси:

\[r_x=r_{2x}-r_{1x}=5-1=4\,\text{м};\]

\[r_y=r_{2y}-r_{1y}=5-2=3\,\text{м}.\]

Тогда перемещение тела за время $t = 2{,}5\,\text{с}$ равно $s=\sqrt{r_x^2+r_y^2}=\sqrt{16+9}=5\,\text{м}.$

Средняя скорость за указанный промежуток времени равна:

\[v _{\text{ср}}= \frac{s}{t} = \frac{5}{2{,}5} =2\,\text{м/с}.\]

Ответ: $2$.
}

% Source sdamgia-11807: https://phys-ege.sdamgia.ru/problem?id=11807
\BankTask{sdamgia-11807}{
Расстояние между пунктами А и В равно $30\,\text{км}$. Из пункта А в направлении пункта В выезжает мотоциклист со скоростью $50\,\text{км/ч}$. Одновременно из пункта В в том же направлении, что и мотоциклист, выезжает трактор со скоростью $20\,\text{км/ч}$. На каком расстоянии от пункта А мотоциклист догонит трактор? (В километрах.)
}{
\smallskip\noindent\textbf{Решение}\par
Для того чтобы догнать трактор, мотоциклисту нужно сократить расстояние в $30\,\text{км}$. Относительная скорость мотоциклиста относительно трактора равна $50\,\text{км/ч} - 20\,\text{км/ч} = 30\,\text{км/ч}$. Следовательно, мотоциклист догонит трактор за $\frac{30\,\text{км}}{30\,\text{км/ч}} = 1\,\text{ч}.$ За это время мотоциклист проедет $50\,\text{км}$.

Ответ: $50$.
}

% Source sdamgia-9752: https://phys-ege.sdamgia.ru/problem?id=9752
\BankTask{sdamgia-9752}{
Скорость течения широкой реки $3{,}6\,\text{км/ч}$. Под каким углом к направлению течения реки лодочник должен направлять лодку, скорость которой относительно воды равна $2\,\text{м/с}$, чтобы за $15$ минут ее снесло по направлению течения на $1{,}8\,\text{км}$?
}{
\smallskip\noindent\textbf{Решение}\par
При движении под углом $\alpha $ к течению реки, лодка будет двигаться по реке со скоростью 

\[V=V_{\text{теч}}+V_{\text{л}}\cos\alpha=3{,}6\cdot\frac{1000}{3600}+2\cos\alpha=(1+2\cos\alpha)\,\text{м/с}.\]

Тогда за время $t$ лодка переместится на расстояние

\[S=Vt= ( 1 + 2 \cdot  \cos  \alpha  ) t,\]

откуда

\[\cos  \alpha  = \frac{\frac{S}{t} - 1}{2} = \frac{1}{2} ( \frac{1800}{900} - 1 ) = \frac{1}{2} .\]

Лодочник должен направлять лодку под углом $60^\circ$ к направлению реки.

Ответ: $60$.
}

% Source sdamgia-23251: https://phys-ege.sdamgia.ru/problem?id=23251
\BankTask{sdamgia-23251}{
Школьник, возвращаясь домой с занятий по подготовке к ЕГЭ по физике, сел на вокзале в стоящую электричку. В это время пошел сильный снег, и вдоль поезда начал дуть ветер с постоянной скоростью. При этом поток падающих за окном снежинок выглядел так, как показано на рис. $1$. Потом поезд поехал, и после его разгона наблюдаемая через окно картина изменилась (см. рис. $2$). Углы наклона к горизонту прямолинейных траекторий снежинок в первом и во втором случаях были равны $\alpha  = 45 ^\circ$ и $\beta  = 10 ^\circ.$ С какой скоростью ехал поезд, если снежинки падают в неподвижном воздухе со скоростью $V = 1\,\text{м/с}$, электричка едет прямолинейно, а скорость и направление ветра не изменяются?
\begin{center}\begin{tikzpicture}[>=stealth,line width=.5pt]

\foreach\shift/\angle/\label in{0/45/1,5/10/2}{\begin{scope}[xshift=\shift cm]\draw(0,0)rectangle(4,3);\foreach\x in{.5,1.2,1.9,2.6,3.3}\foreach\y in{.4,1,1.6,2.2,2.8}\draw(\x,\y)--++(\angle:.15);\draw[thin](1.1,1.5)--(3,1.5);\draw[->,thick](2.6,1.5)--++(180+\angle:1.2);\draw(2.15,1.5)arc(180:180+\angle:.45);\node[below]at(2,-.15){Рис.\ \label};\end{scope}}
\node at(1.9,1.25){$\alpha$};\node at(6.9,1.23){$\beta$};
\end{tikzpicture}\end{center}

}{
\smallskip\noindent\textbf{Решение}\par
$1$. В неподвижной системе отсчета, связанной со стоящим поездом (рис. $1$), скорость снежинок $\overrightarrow V$ по закону сложения скоростей в кинематике складывается из двух составляющих — скорости $\overrightarrow V$ падения снежинки по вертикали и горизонтальной скорости ветра $\vec V_{\text{в}}:\quad\vec V_1=\vec V+\vec V_{\text{в}}.$ Поскольку векторы $\overrightarrow V$ и $\overrightarrow V_{\text{в}}$ взаимно перпендикулярны, а суммарный вектор $\overrightarrow V_1$ наклонен под углом $\alpha  = 45 ^\circ$ к горизонту, то модули $\overrightarrow V$ и $\overrightarrow V_{\text{в}}$ одинаковы и равны $V_{\text{в}}$ = $V = 1\,\text{м/с}$ — ветер дует влево с горизонтальной скоростью $1\,\text{м/с}$.

$2$. В системе отсчета, движущейся вместе с поездом (рис. $2$), снежинки приобретают относительно него дополнительную горизонтальную скорость $V_{\text{п}}$, направленную также влево (против вектора скорости поезда $\overrightarrow V_{\text{п}}$) и складывающуюся со скоростью ветра.

$3$. При этом вертикальная составляющая скорости снежинок остается равной $V = 1\,\text{м/с}$. Таким образом, как следует из прямоугольного треугольника для сложения горизонтальной и вертикальной составляющих скорости снежинки относительно поезда, $V = ( V_{\text{п}} + V_{\text{в}} ) \cdot  \tan  \beta  ,$ откуда $V_{\text{п}}=V/\tan\beta-V_{\text{в}}\approx1/0{,}176-1\approx4{,}67\,\text{м/с}\approx17\,\text{км/ч}.$

Ответ: $V_{\text{п}}\approx17\,\text{км/ч}.$
}


% Source 46782: https://phys-ege.sdamgia.ru/problem?id=46782
\BankTask{46782}{
Камень, поднятый над поверхностью земли, падает без начальной скорости вертикально вниз и за последнюю секунду своего движения проходит путь $\Delta h=10\,\text{м}$. Определите среднюю скорость камня за всё время падения. Сопротивлением воздуха пренебречь.
}{
\smallskip\noindent\textbf{Решение}\par
Средняя скорость камня равна отношению высоты H, с которой бросили камень, к времени его падения $t_0$ до поверхности земли: $v_{\text{ср}} = \frac{H}{t_0} .$
Камень движется вертикально вниз равноускоренно с ускорением свободного падения, следовательно: $H = \frac{g t_0 ^2}{2} ,$ откуда $v_{\text{ср}} = \frac{g t_0}{2} .$ 
Найдём путь $\Delta h,$ пройденный камнем за последнюю секунду падения: 

\[\Delta h = H - h = \frac{g t_0 ^2}{2} - \frac{g ( t_0 - \Delta t ) ^2}{2} ,\]

где $\Delta t = 1 \,\text{с}$ — время, за которое камень прошел путь $\Delta h.$ Отсюда получим: 

$t_0 = \frac{2 \Delta h + g \Delta t ^2}{2 g \Delta t} = \frac{2 \cdot 10 + 10 \cdot 1 ^2}{2 \cdot 10 \cdot 1} = 1{,}5 \,\text{с},$ 
$v_{\text{ср}} = \frac{g t_0}{2} = \frac{10 \cdot 1{,}5}{2} = 7{,}5 \,\text{м/с}.$

Ответ: $v _{\,\text{ср}} = 7{,}5 \,\text{м/с}.$
}

\end{document}
